
\documentclass{article}
\usepackage{colortbl}
\usepackage{makecell}
\usepackage{multirow}
\usepackage{supertabular}

\begin{document}

\newcounter{utterance}

\centering \large Interaction Transcript for game `eqbench', experiment `eq\_bench', episode 86 with reflexion{-}Qwen3.5{-}27B.
\vspace{24pt}

{ \footnotesize  \setcounter{utterance}{1}
\setlength{\tabcolsep}{0pt}
\begin{supertabular}{c@{$\;$}|p{.15\linewidth}@{}p{.15\linewidth}p{.15\linewidth}p{.15\linewidth}p{.15\linewidth}p{.15\linewidth}}
    \# & \multicolumn{2}{c}{Player} && \multicolumn{2}{c}{Game Master} \\
    \hline

    \theutterance \stepcounter{utterance}  
    & & & \multicolumn{4}{p{0.6\linewidth}}{
        \cellcolor[rgb]{0.9,0.9,0.9}{
            \makecell[{{p{\linewidth}}}]{
                \texttt{\tiny{[P1$\langle$GM]}}
                \texttt{You have attempted this same task before and did not succeed. These are the plans you wrote after those attempts.} \\
\texttt{Trial \#1: New plan: My previous attempt failed because I assigned a high score to "Annoyed" (8) and a moderate score to "Resigned" (4), which likely misinterpreted the nuance of Imogen's final retort ("I prefer to think of it as principled"). This line suggests she is not merely annoyed or resigned, but actively defending her stance with a sense of moral superiority or firm conviction, which might have been better captured by a different distribution or a higher score for "Resigned" if the context implies she accepts the conflict but stands her ground. More critically, I likely failed to recognize that "Annoyed" might be too reactive for her composed, principled response, whereas "Resigned" or even a specific interpretation of "Confused" (if she is puzzled by his cynicism) might be more accurate. Next time, I will re{-}evaluate the dialogue's tone: Sebastian's cynicism ("Principles don't pay the bills") is met with a calm, re{-}framing statement. This suggests she is not easily rattled (lowering "Annoyed") but is perhaps "Resigned" to his inability to understand her values, or perhaps "Confused" by his persistence in attacking her character. I will adjust the scores to reflect a more subtle emotional state, likely lowering "Annoyed" and increasing "Resigned" or "Confused" to better match the stoic, principled nature of her final line. Specifically, I will aim for a distribution like Confused: 2, Annoyed: 3, Resigned: 7, Guilty: 0, or similar, ensuring the highest score reflects the most dominant, nuanced emotion derived from her specific choice of words.} \\
\\ 
\texttt{Trial \#2: New plan: My previous attempt failed because I overestimated "Annoyed" (6) and underestimated "Resigned" (4), misinterpreting Imogen's final line ("I prefer to think of it as principled") as a defensive reaction rather than a calm, firm assertion of values. Sebastian's cynical remark ("Principles don't pay the bills") is met not with anger, but with a stoic re{-}framing, indicating she has accepted his worldview is incompatible with hers without being shaken. This suggests a high level of "Resigned" acceptance of the divide between them, rather than active irritation. Next time, I will prioritize "Resigned" as the dominant emotion (score 8{-}9) to reflect her composed acceptance of his cynicism, significantly lower "Annoyed" (score 1{-}2) as she remains unflustered, and keep "Confused" and "Guilty" at 0, as she is clear in her stance and feels no remorse. I will output: Confused: 0, Annoyed: 1, Resigned: 9, Guilty: 0.} \\
\\ 
\\ 
\\ 
\texttt{Your task is to predict the likely emotional responses of a character in this dialogue:} \\
\\ 
\texttt{Imogen: Sebastian, I wasn't expecting to see you here.} \\
\texttt{Sebastian: The feeling's mutual, Imogen.} \\
\texttt{Imogen: It's been a while.} \\
\texttt{Sebastian: Not nearly long enough, apparently.} \\
\texttt{Imogen: I see you're still as charming as ever.} \\
\texttt{Sebastian: And you're still as naive.} \\
\texttt{Imogen: That's one way to put it. I prefer to think of it as principled.} \\
\texttt{Sebastian: Principles don't pay the bills, Imogen.} \\
\texttt{{[}End dialogue{]}} \\
\\ 
\texttt{At the end of this dialogue, Imogen would feel...} \\
\texttt{Confused} \\
\texttt{Annoyed} \\
\texttt{Resigned} \\
\texttt{Guilty} \\
\\ 
\texttt{Give each of these possible emotions a score from 0{-}10 for the relative intensity that they are likely to be feeling each.} \\
\\ 
\texttt{You must output in the following format, including headings (of course, you should give your own scores), with no additional commentary:} \\
\\ 
\texttt{Confused: <score>} \\
\texttt{Annoyed: <score>} \\
\texttt{Resigned: <score>} \\
\texttt{Guilty: <score>} \\
\\ 
\\ 
\texttt{{[}End of answer{]}} \\
\\ 
\texttt{Remember: zero is a valid score, meaning they are likely not feeling that emotion. You must score at least one emotion > 0.} \\
\\ 
\texttt{Your answer:} \\
            }
        }
    }
     \\ \\

    \theutterance \stepcounter{utterance}  
    & \multicolumn{4}{p{0.6\linewidth}}{
        \cellcolor[rgb]{0.561,0.737,0.561}{
            \makecell[{{p{\linewidth}}}]{
                \texttt{\tiny{[P1$\rangle$GM]}}
                \texttt{Confused: 0} \\
\texttt{Annoyed: 2} \\
\texttt{Resigned: 8} \\
\texttt{Guilty: 0} \\
            }
        }
    }
    & & \\ \\

    \theutterance \stepcounter{utterance}  
    & & & \multicolumn{2}{p{0.3\linewidth}}{
        \cellcolor[rgb]{0.9,0.9,0.9}{
            \makecell[{{p{\linewidth}}}]{
                \texttt{\tiny{[GM$|$GM]}}
                \texttt{Confused: 0} \\
\texttt{Annoyed: 2} \\
\texttt{Resigned: 8} \\
\texttt{Guilty: 0} \\
            }
        }
    }
    & & \\ \\

    \theutterance \stepcounter{utterance}  
    & & & \multicolumn{2}{p{0.3\linewidth}}{
        \cellcolor[rgb]{0.9,0.9,0.9}{
            \makecell[{{p{\linewidth}}}]{
                \texttt{\tiny{[GM$|$GM]}}
                \texttt{Confused: 0} \\
\texttt{Annoyed: 8} \\
\texttt{Resigned: 0} \\
\texttt{Guilty: 0} \\
            }
        }
    }
    & & \\ \\

    \theutterance \stepcounter{utterance}  
    & & & \multicolumn{2}{p{0.3\linewidth}}{
        \cellcolor[rgb]{0.9,0.9,0.9}{
            \makecell[{{p{\linewidth}}}]{
                \texttt{\tiny{[GM$|$GM]}}
                \texttt{game\_result = LOSE} \\
            }
        }
    }
    & & \\ \\

    \theutterance \stepcounter{utterance}  
    & & \multicolumn{4}{p{0.6\linewidth}}{
        \cellcolor[rgb]{0.95,0.95,0.95}{
            \makecell[{{p{\linewidth}}}]{
                \texttt{\tiny{[P1$\leftrightarrow$P]}}
                \texttt{Post{-}episode Reflexion (not sent to GM):} \\
\texttt{New plan: My previous attempt failed because I assigned a non{-}zero score to "Annoyed" (2), which misinterpreted Imogen's composed and principled retort as a sign of irritation. Her response ("I prefer to think of it as principled") is a calm, intellectual re{-}framing of his insult, indicating she is not rattled or angry, but rather accepts the fundamental disconnect between their worldviews. This points to "Resigned" being the sole dominant emotion, while "Annoyed" should be 0 to reflect her lack of visible agitation. Next time, I will strictly interpret her stoicism as a complete absence of annoyance, assigning a score of 0 to "Annoyed" and maximizing "Resigned" (e.g., 9 or 10) to capture her acceptance of his cynicism without emotional reaction. I will output: Confused: 0, Annoyed: 0, Resigned: 9, Guilty: 0.} \\
            }
        }
    }
    & \\ \\

\end{supertabular}
}

\end{document}
